% ============================================================
\chapter{Convergence Spaces and Filters}
\label{ch:convergence}
% ============================================================

What is convergence, really? In a standard analysis course, we learn that a sequence $(x_n)$ converges to $\ell$ if, for every neighbourhood of $\ell$, the sequence eventually stays inside it. But this definition presupposes a topology. In the 1960s, mathematicians such as Gustave Choquet, H.\,J.\,Kowalsky, and C.\,H.\,Cook asked a bold question: what if convergence itself were the primitive notion? Rather than defining open sets first and deducing convergence, why not specify directly which filters converge to which points? This logical inversion gave birth to \emph{convergence spaces}, a generalization of topological spaces that elegantly solves the cartesian closure problem --- a technical defect of the category of topological spaces that complicated the study of function spaces.

\begin{intuition}
Convergence in a topological space is classically described via filters or nets.
Convergence spaces generalize topological spaces by directly specifying which filters
converge to which points, without requiring the convergence structure to arise from a
topology. This generalization solves the cartesian closure problem and is particularly
suited to function spaces.
\end{intuition}

% ------------------------------------------------------------
\section{Filters: Review and Complements}
% ------------------------------------------------------------

\begin{definition}[Filter]
A \emph{filter} on a set $X$ is a family $\mathcal{F} \subseteq \mathcal{P}(X)$ such that:
\begin{enumerate}
  \item $X \in \mathcal{F}$ and $\emptyset \notin \mathcal{F}$;
  \item if $A, B \in \mathcal{F}$, then $A \cap B \in \mathcal{F}$;
  \item if $A \in \mathcal{F}$ and $A \subseteq B$, then $B \in \mathcal{F}$.
\end{enumerate}
\end{definition}

\begin{definition}[Filter base]
A \emph{filter base} is a nonempty family $\mathcal{B} \subseteq \mathcal{P}(X)$, not
containing $\emptyset$, such that for all $B_1, B_2 \in \mathcal{B}$, there exists
$B_3 \in \mathcal{B}$ with $B_3 \subseteq B_1 \cap B_2$. The generated filter is
$\mathcal{F} = \{A \subseteq X : \exists B \in \mathcal{B},\, B \subseteq A\}$.
\end{definition}

\begin{definition}[Ultrafilter]
An \emph{ultrafilter} is a maximal filter (under inclusion). Equivalently, $\mathcal{U}$
is an ultrafilter if for every $A \subseteq X$, either $A \in \mathcal{U}$ or
$X \setminus A \in \mathcal{U}$.
\end{definition}

\begin{theorem}[Existence of ultrafilters]
Every filter is contained in an ultrafilter (by Zorn's lemma).
\end{theorem}

\begin{definition}[Image filter]
If $f\colon X \to Y$ and $\mathcal{F}$ is a filter on $X$, the image filter is
$f(\mathcal{F}) = \{B \subseteq Y : f^{-1}(B) \in \mathcal{F}\}$.
\end{definition}

% ------------------------------------------------------------
\section{Convergence in Topological Spaces}
% ------------------------------------------------------------

\begin{definition}[Neighborhood filter]
For $x \in X$, the \emph{neighborhood filter} $\mathcal{V}(x)$ is the set of neighborhoods
of $x$: $\mathcal{V}(x) = \{V \subseteq X : \exists U \in \tau,\, x \in U \subseteq V\}$.
\end{definition}

\begin{definition}[Filter convergence]
A filter $\mathcal{F}$ on a topological space $X$ \emph{converges} to $x$ (written
$\mathcal{F} \to x$) if $\mathcal{V}(x) \subseteq \mathcal{F}$.
\end{definition}

\begin{theorem}[Characterization of topology via filters]
Let $X$ be a topological space. Then:
\begin{enumerate}
  \item $U$ is open iff for every $x \in U$ and every filter $\mathcal{F} \to x$,
    $U \in \mathcal{F}$.
  \item $f\colon X \to Y$ is continuous iff for every filter $\mathcal{F} \to x$,
    $f(\mathcal{F}) \to f(x)$.
  \item $X$ is compact iff every ultrafilter converges.
  \item $X$ is Hausdorff iff every filter has at most one limit.
\end{enumerate}
\end{theorem}

\begin{proof}
(3) $(\Rightarrow)$ Let $\mathcal{U}$ be an ultrafilter. For each $x \in X$, if
$\mathcal{U} \not\to x$, there exists $U_x \in \mathcal{V}(x)$ with
$U_x \notin \mathcal{U}$, so $X \setminus U_x \in \mathcal{U}$. If no ultrafilter
converges, the $U_x$ cover $X$; extract a finite subcover $U_{x_1}, \ldots, U_{x_n}$;
then $\bigcap_{i=1}^n (X \setminus U_{x_i}) = \emptyset \in \mathcal{U}$, contradiction.

$(\Leftarrow)$ Let $(U_i)$ be an open cover with no finite subcover. The sets
$X \setminus U_i$ have the finite intersection property, so they generate a filter
contained in an ultrafilter $\mathcal{U}$. By hypothesis, $\mathcal{U} \to x$ for some
$x \in U_j$, giving $U_j \in \mathcal{U}$ and $X \setminus U_j \in \mathcal{U}$,
contradiction.
\end{proof}

% ------------------------------------------------------------
\section{Convergence Spaces}
% ------------------------------------------------------------

\begin{definition}[Convergence space]
A \emph{convergence space} is a pair $(X, \to)$ where $X$ is a set and $\to$ is a relation
between filters on $X$ and points of $X$ (called a \emph{convergence structure})
satisfying:
\begin{enumerate}
  \item For every $x \in X$, the principal filter
    $\dot{x} = \{A \subseteq X : x \in A\}$ converges to $x$: $\dot{x} \to x$.
  \item If $\mathcal{F} \to x$ and $\mathcal{F} \subseteq \mathcal{G}$, then
    $\mathcal{G} \to x$.
\end{enumerate}
\end{definition}

\begin{definition}[Pretopological convergence space]
A convergence space is \emph{pretopological} if additionally:
\begin{enumerate}
  \setcounter{enumi}{2}
  \item If $\mathcal{F} \to x$ and $\mathcal{G} \to x$, then
    $\mathcal{F} \cap \mathcal{G} \to x$.
\end{enumerate}
\end{definition}

\begin{definition}[Topological convergence space]
A convergence space is \emph{topological} if additionally the neighborhood filter
$\mathcal{V}(x) = \bigcap\{\mathcal{F} : \mathcal{F} \to x\}$ converges to $x$.
\end{definition}

\begin{theorem}[Hierarchy]
We have strict inclusions of categories:
\[
  \mathbf{Top} \subsetneq \mathbf{PrTop} \subsetneq \mathbf{Conv}
\]
where $\mathbf{Conv}$ is the category of convergence spaces and continuous maps
(preserving convergence), and $\mathbf{PrTop}$ is the pretopological subcategory.
\end{theorem}

% ------------------------------------------------------------
\section{Cartesian Closedness}
% ------------------------------------------------------------

\begin{theorem}[$\mathbf{Conv}$ is cartesian closed]
The category $\mathbf{Conv}$ is cartesian closed: for all convergence spaces $X$ and $Y$,
the space $[X,Y]$ of continuous maps with the continuous convergence structure is a
convergence space, and:
\[
  \mathrm{Hom}_{\mathbf{Conv}}(X \times Y, Z) \cong
  \mathrm{Hom}_{\mathbf{Conv}}(X, [Y, Z]).
\]
\end{theorem}

\begin{definition}[Continuous convergence]
On $[X,Y]$, a filter $\Phi$ \emph{converges continuously} to $f$ if for every $x \in X$
and every filter $\mathcal{F} \to x$ in $X$, the evaluation filter $\Phi[\mathcal{F}]$
(generated by $\{g(y) : g \in \phi, y \in F\}$ for $\phi \in \Phi$,
$F \in \mathcal{F}$) converges to $f(x)$ in $Y$.
\end{definition}

\begin{remark}
The fact that $\mathbf{Conv}$ is cartesian closed while $\mathbf{Top}$ is not is one of
the main motivations for studying convergence spaces. In denotational semantics, this
allows constructing function spaces without leaving the category.
\end{remark}

% ------------------------------------------------------------
\section{Nets and Convergence}
% ------------------------------------------------------------

\begin{definition}[Net]
A \emph{net} in $X$ is a function $\varphi\colon I \to X$ where $(I, \leq)$ is a directed
set. We write $(x_i)_{i\in I}$.
\end{definition}

\begin{definition}[Net convergence]
A net $(x_i)_{i\in I}$ converges to $x$ in a topological space if for every neighborhood
$V$ of $x$, there exists $i_0 \in I$ such that for all $i \geq i_0$, $x_i \in V$.
\end{definition}

\begin{theorem}[Filter-net equivalence]
In a topological space, a set $F$ is closed iff every limit of a net in $F$ lies in $F$.
Similarly, $f$ is continuous iff it preserves net limits.
\end{theorem}

\begin{warning}
In non-Hausdorff spaces, a net (or filter) may converge to multiple points simultaneously.
The set of limits of a filter $\mathcal{F}$ in a $T_0$-space is a downward-closed set for
the specialization order.
\end{warning}

% ------------------------------------------------------------
\section{Scott Filters and Convergence}
% ------------------------------------------------------------

\begin{definition}[Scott filter]
Let $D$ be a dcpo. A filter $\mathcal{F}$ on $D$ is a \emph{Scott filter} if it contains a
directed set $S$ such that $\bigsqcup S$ exists and $\mathcal{F}$ converges to
$\bigsqcup S$ in the Scott topology.
\end{definition}

\begin{proposition}
In a dcpo $D$ with the Scott topology, a filter $\mathcal{F}$ converges to $x$ if and only
if every Scott-open set containing $x$ belongs to $\mathcal{F}$.
\end{proposition}

% ------------------------------------------------------------
\section{Choquet Spaces}
% ------------------------------------------------------------

\begin{definition}[Choquet space]
A topological space $X$ is a \emph{Choquet space} (or has a favorable topology) if Player~II
has a winning strategy in the Choquet game: the players alternately choose nonempty opens
$U_1 \supseteq V_1 \supseteq U_2 \supseteq V_2 \supseteq \ldots$, and Player~II wins if
$\bigcap_n V_n \neq \emptyset$.
\end{definition}

\begin{theorem}
Every complete metric space is a Choquet space. Every continuous dcpo with the Scott
topology is a Choquet space.
\end{theorem}

% ------------------------------------------------------------
\section{Convergence and Domain Theory}
% ------------------------------------------------------------

\begin{proposition}[Scott convergence and approximation]
In a continuous domain $D$, a filter $\mathcal{F}$ converges to $x$ in the Scott topology
iff for every $b \ll x$, $\uparrow\! b \in \mathcal{F}$. Convergence in a continuous
domain is thus determined by the way-below relation $\ll$.
\end{proposition}

\begin{theorem}[Convergence spaces and domains]
The category of continuous domains with Scott-continuous functions is equivalent to a full
subcategory of $\mathbf{Conv}$ consisting of convergence spaces satisfying additional
axioms (directed convergence).
\end{theorem}

% ------------------------------------------------------------
\section{Exercises}
% ------------------------------------------------------------

\begin{exercise}
Show that a topological space is $T_0$ iff in the associated convergence structure, the
set of limits of every filter is a downward-closed set (for the specialization order) and
two distinct points do not have the same convergent filters.
\end{exercise}

\begin{exercise}
Construct a convergence space that is not pretopological.
\end{exercise}

\begin{exercise}
Show that continuous convergence on $[X,Y]$ coincides with compact-open convergence when
$X$ is locally compact Hausdorff.
\end{exercise}

\begin{exercise}
Let $D$ be a dcpo. Show that the filter generated by a directed set $S$ converges to
$\bigsqcup S$ in the Scott topology.
\end{exercise}

\begin{exercise}
Show that in a $T_0$-space, if a filter $\mathcal{F}$ converges to $x$ and to $y$ with
$x \leq y$, then $\mathcal{F}$ also converges to every $z$ with $x \leq z \leq y$.
\end{exercise}

\begin{exercise}[Category $\mathbf{Conv}$ and limits]
Show that $\mathbf{Conv}$ is complete and cocomplete. Describe initial and final structures
in terms of convergence.
\end{exercise}
